%% kapa: exam=1 round=30 year=2019 period=1 form=A subject=civil_law subject_ko=민법 seq=14 number=14
%% source: Q-Net 감정평가사 기출문제 (cst00309) artlSeq=5205205 file=extracted/1008/5205205/2019년 제30회 감정평가사 1차시험 문제지/2019년 제30회 감정평가사 1차시험 문제지/2019년도 제30회 감정평가사 1차 1교시 A형.hwp
%% conversion: kordoc 4.24.51 -> tools/gen1.py (auto); review=proofread-sonnet; flags=none
\begin{question}{14}
대리에 관한 설명으로 옳지 않은 것은? (다툼이 있으면 판례에 따름)
\begin{choices}
  \item 불법행위에는 대리의 법리가 적용되지 않는다.
  \item 대리인이 자신의 이익을 도모하기 위하여 대리권을 남용한 경우는 무권대리에 해당한다.
  \item 대리인의 대리행위가 공서양속에 반하는 경우, 본인이 그 사정을 몰랐다고 하더라도 그 행위는 무효이다.
  \item 대리인이 상대방에게 사기ㆍ강박을 하였다면 상대방은 본인이 그에 대해 선의․무과실이라 하더라도 대리인과 행한 법률행위를 취소할 수 있다.
  \item 복대리인은 본인의 대리인이다.
\end{choices}
\end{question}
\answer{2}
